\section{Introduction}
\label{sec:intro}
\let\thefootnote\relax\footnotetext{$^*$Lead authors, ordered alphabetically, see \cref{sec:contributions} for list of contributions. \\Correspondence to \href{mailto:dibya.ghosh@berkeley.edu,homer_walke@berkeley.edu,pertsch@berkeley.edu,kvablack@berkeley.edu,oier.mees@eecs.berkeley.edu}{\scriptsize\texttt{\{dibya.ghosh, homer\_walke, pertsch, kvablack, oier.mees\}@berkeley.edu}}}
\setcounter{footnote}{0}
The common approach for robotic learning is to train policies on datasets collected for the specific robot and task at hand. Learning from scratch in this way requires significant data collection effort for each task, and the resulting policies usually exhibit only narrow generalization. In principle, collected experience from other robots and tasks offers a possible solution, exposing models to a diverse set of robotic control problems that may improve generalization and performance on downstream tasks.
However, even as general-purpose models become ubiquitous in natural language \citep{openaiGPT4TechnicalReport2023,touvronLLaMAOpenEfficient2023}) and computer vision \citep{rombachHighResolutionImageSynthesis2022, kirillovSegmentAnything2023}, it has proven challenging to build the analogous ``general-purpose robot model'' that can control many robots for many tasks. Training a unified control policy in robotics presents unique challenges, requiring handling different robot embodiments, sensor setups, action spaces, task specifications, environments, and compute budgets.

Towards this direction, several works have proposed robotic foundation models that directly map robot observations to actions and provide zero-shot or few-shot generalization to new domains and robots.
We broadly refer to these models as ``generalist robot policies'' (GRPs), 
emphasizing their ability to perform low-level visuomotor control across tasks, environments, and robotic systems~\citep{reed2022generalist,bousmalis2023robocat,driess2023palm,zitkovich2023rt,brohan2022rt,shah2023vint,scaleblog,wayveblog,hu2023gaia1,yang2023polybot,kumar2023robohive}. For example, the GNM model~\citep{shahGNMGeneralNavigation2023} generalizes across different robotic navigation scenarios, the RoboCat model~\citep{bousmalis2023robocat} handles different robot embodiments for goal-conditioned tasks, and the RT-X model~\citep{open_x_embodiment_rt_x_2023} performs language-conditioned manipulation across five robot embodiments. 
Although these models represent significant steps toward a true ``general-purpose robot model,'' they have been limited in multiple important aspects: they typically constrain downstream users to a pre-defined and often restrictive set of input observations, \eg a single camera stream; they lack support for effective finetuning to new domains; and importantly, the largest of these models are not available to the general public. 

We design a system for pretraining generalist robot policies more suitable for the diversity of interfaces in downstream robotic applications. The core of our model is a transformer architecture that maps arbitrary input tokens (created from observations and tasks) to output tokens (then decoded into actions), which can be trained on a diverse dataset of robots and tasks. With no additional training, this policy can accept different camera configurations (e.g., workspace or wrist cameras), can control different robots, and can be guided via either language commands or goal images --- all by simply changing which tokens are fed into the model. Most importantly, the model can be adapted to new robot setups with new sensory inputs, action spaces, or morphologies by adding appropriate adapters and finetuning with a small target domain dataset and an accessible compute budget.

Our primary contribution is \methodfancy{}, a transformer-based policy pretrained on the largest robot manipulation dataset to date: \ntrajs{} robot demonstrations from the Open X-Embodiment dataset~\citep{open_x_embodiment_rt_x_2023}. \method{} is the first GRP that can be effectively finetuned to new observations and action spaces and the first generalist robot manipulation policy that is fully open-source, including the training pipeline, model checkpoints, and data. Finally, while the individual components that comprise \method{} --- a transformer backbone, support for both language and goal image specification, and a diffusion head to model expressive action distributions --- have been discussed in prior work, the particular combination of these components into a powerful generalist robot policy is unique and novel. 


We demonstrate through extensive experiments on \nsetups{}~robots across \ninstitutions{}~institutions that our combined system leads to state-of-the-art performance for out-of-the-box multi-robot control for single and dual-arm manipulation tasks and that \method{} can be used as an effective initialization for finetuning to unseen setups with new observation and action spaces. In the process, we carefully study the effect of different design decisions when pretraining GRPs; we evaluate how the choice of data distribution, model architecture, and policy formulation affects the quality of the pretrained GRP. Our evaluation highlights the utility of scale and flexibility: our best models are those trained on the widest data mixtures, with the least restrictive inductive biases, and with policy objectives that can fit the diversity of behaviors in the pretraining data.


Along with this paper, we release all resources required to train, use, reproduce, and finetune an \method{} model. We provide pretrained \method{} model checkpoints with \smallnparams{} and \largenparams{} parameters that, out of the box, support multiple RGB camera inputs as well as both language and goal image task specification. We also provide scripts for finetuning these models on new domains, as well as our complete pretraining pipeline, including optimized data loaders, transformer implementations for multimodal inputs, and tools to monitor training progress.



